\documentclass[10pt,a4paper]{amsart}
\usepackage[margin=19mm]{geometry}
\usepackage{amsmath,amssymb,mathtools}
\usepackage[scr=rsfs,cal=euler]{mathalfa}
\usepackage{xfrac}
\usepackage[unicode,hidelinks,bookmarks=false]{hyperref}
\newcommand{\ie}{i.e.\ }
\newcommand{\cf}{cf.\ }
\newcommand{\resp}{respectively\ }
\newcommand{\Subsection}{\subsection}
\providecommand{\nobreakdash}{\nobreak\hbox{-}}
\setlength{\emergencystretch}{2em}
\multlinegap0pt
\usepackage[shortlabels]{enumitem}
\usepackage{amssymb}
\usepackage{mathtools}
\usepackage{bm}
\newtheorem{assumption}{Assumption}
\newtheorem{lemma}{Lemma}
\newcommand{\dist}{\mathop{\mathrm{dist}}}
\newcommand{\tr}{\mathop{\mathrm{tr}}}
\title{Consistency of the Bayesian Information Criterion for Model Selection in Exploratory Factor Analysis}
\author{Hien Duy Nguyen, Kei Hirose}
\date{Source: arXiv:2604.07998v1 (CC BY 4.0)}
\begin{document}
\maketitle

\noindent\emph{Source and licence.} Consistency of the Bayesian Information Criterion for Model Selection in Exploratory Factor Analysis. Version arXiv:2604.07998v1 (CC BY 4.0), distributed by arXiv under the Creative Commons Attribution 4.0 International licence, http://creativecommons.org/licenses/by/4.0/, as recorded in the arXiv licence metadata for this exact version and shown on the abstract page https://arxiv.org/abs/2604.07998v1. Full licence notice: \texttt{licenses/arxiv-2604.07998-CC-BY-4.0.txt}.\par\smallskip

\noindent\textit{Editorial scope.} Counted unit: the article's lemma lem:exactoverfit (source0.tex lines 1312--1323). The article's complete original proof of it is delivered with it (source0.tex lines 1324--1388, one \texttt{proof} environment of 65 lines). No mathematics is written, repaired or re-derived: the statement and the proof are the article's own lines, in the article's own notation, and no retained line is substituted or re-flowed.  Retained uncounted as the source-local context the proof needs: lines 443--458; lines 1265--1279.  Nothing was omitted from any retained span, so the package carries no omission record.  No retained line contains a citation, so the wrapper carries no bibliography and no bibliography entry.  No retained line contains a figure environment, an image-inclusion command or drawing code, so no image is delivered and the package declares no asset dependency.  Editorial formatting only, in addition: an AMS \texttt{amsart} preamble that restates the article's own declarations and macros used by the retained lines, namely the environments \texttt{assumption}, \texttt{lemma} on separate counters, together with the standard packages those lines need.  The retained environments print in this document as Assumption 1, Lemma 1 in order of appearance, whereas the article prints the same items with the numbering of its own full document.  The complete native source of the article as distributed in the arXiv e-print for this exact version is bundled unchanged as \texttt{sources/198127/source0.tex}.
\begin{assumption}\label{ass:miss} \leavevmode 
\begin{enumerate}[label=\textnormal(M\arabic*)]
\item \label{it:M1} The observations are iid with mean $\mu_{0}$, covariance
$\Sigma_{0}\in\mathbb{S}_{++}^{p}$, and finite fourth moment $\mathbb{E}\|X_{1}\|^{4}<\infty$. 
\item \label{it:M2} There exists a nonempty compact set $\mathcal{G}_{*}\subset\mathcal{K}$
such that 
\[
\mathcal{G}_{k}=\mathcal{G}_{*}\qquad\text{for every }k\in K_{*}.
\]
\item \label{it:M3} For each $k\in K_{*}$ there exist constants $c_{k}>0$
and $\eta_{k}>0$ such that 
\[
Q(\Sigma)\le V_{*}-c_{k}\,\dist(\Sigma,\mathcal{G}_{*})^{2}\qquad\text{for all }\Sigma\in\mathcal{F}_{k}\text{ with }\dist(\Sigma,\mathcal{G}_{*})<\eta_{k}.
\]
\end{enumerate}
\end{assumption}

\begin{lemma}[Uniform approximation on compact covariance sets]\label{lem:uniform}
Let $\mathcal{S}\subset\mathbb{S}_{++}^{p}$ be compact. Then there
exists a finite constant $C_{\mathcal{S}}$ such that 
\[
\sup_{\Sigma\in\mathcal{S}}\left|\widetilde{\ell}_{n}(\Sigma)-nQ(\Sigma)\right|\le C_{\mathcal{S}}n\|S_{n}-\Sigma_{0}\|_{F}.
\]
Under Assumption~\ref{it:M1}, 
\[
\sup_{\Sigma\in\mathcal{S}}\left|\widetilde{\ell}_{n}(\Sigma)-nQ(\Sigma)\right|=O\left(\sqrt{n\log\log n}\right)\qquad\text{a.s.}
\]
In particular, 
\[
\sup_{\Sigma\in\mathcal{S}}\left|\frac{1}{n}\widetilde{\ell}_{n}(\Sigma)-Q(\Sigma)\right|\to0\qquad\text{a.s.}
\]
\end{lemma}

\begin{lemma}[Likelihood gain from overfitting]\label{lem:exactoverfit}
Take Assumption~\ref{ass:miss} and fix $k\in K_{*}$. Then there
exists an almost surely finite random constant $C_{k}$ such that
\[
T_{k,n}\le U_{n}+C_{k}\log\log n
\]
for all sufficiently large $n$, almost surely. Since $\mathcal{G}_{*}\subset\mathcal{F}_{k}$,
one also has $T_{k,n}\ge U_{n}$, and thus 
\[
T_{k,n}-U_{n}=O(\log\log n)\qquad\text{a.s.}
\]
\end{lemma}

\begin{proof}
Fix $k\in K_{*}$. For $\Sigma\in\mathcal{F}_{k}$, write $d(\Sigma)=\dist(\Sigma,\mathcal{G}_{*})$
and, because $\mathcal{G}_{*}$ is compact, choose a projection point
$\Pi(\Sigma)\in\mathcal{G}_{*}$ with $\|\Sigma-\Pi(\Sigma)\|_{F}=d(\Sigma)$.
Let $\Delta_{n}=S_{n}-\Sigma_{0}$. Since $U_{n}\ge\widetilde{\ell}_{n}(\Pi(\Sigma))$, 
\[
\widetilde{\ell}_{n}(\Sigma)-U_{n}\le\widetilde{\ell}_{n}(\Sigma)-\widetilde{\ell}_{n}(\Pi(\Sigma)).
\]
Using the identity from Lemma~\ref{lem:uniform}, 
\[
\widetilde{\ell}_{n}(\Sigma)-\widetilde{\ell}_{n}(\Pi(\Sigma))=n\left(Q(\Sigma)-V_{*}\right)-\frac{n}{2}\tr\left(\Delta_{n}(\Sigma^{-1}-\Pi(\Sigma)^{-1})\right).
\]
The inverse map is Lipschitz on $\mathcal{K}$, so 
\[
\|\Sigma^{-1}-\Pi(\Sigma)^{-1}\|_{F}\le\underline{\psi}^{-2}d(\Sigma).
\]
Hence 
\[
\widetilde{\ell}_{n}(\Sigma)-U_{n}\le n\left(Q(\Sigma)-V_{*}\right)+A_{1}\sqrt{n\log\log n}\,d(\Sigma)
\]
eventually almost surely, for some almost surely finite random constant
$A_{1}$.

Now split $\mathcal{F}_{k}$ into a near region $\{d(\Sigma)<\eta_{k}\}$
and a far region $\{d(\Sigma)\ge\eta_{k}\}$. On the near region,
Assumption~\ref{it:M3} implies 
\[
Q(\Sigma)-V_{*}\le-c_{k}d(\Sigma)^{2},
\]
and therefore 
\begin{align*}
\widetilde{\ell}_{n}(\Sigma)-U_{n} & \le-nc_{k}d(\Sigma)^{2}+A_{1}\sqrt{n\log\log n}\,d(\Sigma)\\
 & =-nc_{k}\left(d(\Sigma)-\frac{A_{1}}{2c_{k}}\sqrt{\frac{\log\log n}{n}}\right)^{2}+\frac{A_{1}^{2}}{4c_{k}}\log\log n\\
 & \le\frac{A_{1}^{2}}{4c_{k}}\log\log n.
\end{align*}
On the far region, set 
\[
\mathcal{F}_{k}^{\mathrm{far}}=\{\Sigma\in\mathcal{F}_{k}:d(\Sigma)\ge\eta_{k}\}.
\]
If $\mathcal{F}_{k}^{\mathrm{far}}=\varnothing$, then the near-region
bound already proves the claim. Otherwise, because $d(\cdot)$
is continuous and $\mathcal{F}_{k}$ is compact, the set $\mathcal{F}_{k}^{\mathrm{far}}$
is compact. If $\sup_{\Sigma\in\mathcal{F}_{k}^{\mathrm{far}}}Q(\Sigma)=V_{*}$,
then continuity of $Q$ on $\mathcal{F}_{k}^{\mathrm{far}}$ would
yield some $\Sigma^{\dagger}\in\mathcal{F}_{k}^{\mathrm{far}}$ with
$Q(\Sigma^{\dagger})=V_{*}$. But then $\Sigma^{\dagger}\in\mathcal{G}_{k}=\mathcal{G}_{*}$
by Assumption~\ref{it:M2}, which contradicts $d(\Sigma^{\dagger})\ge\eta_{k}>0$.
Therefore 
\[
b_{k}=V_{*}-\sup_{\Sigma\in\mathcal{F}_{k}^{\mathrm{far}}}Q(\Sigma)>0,
\]
so $Q(\Sigma)\le V_{*}-b_{k}$ whenever $d(\Sigma)\ge\eta_{k}$. Hence
\[
\widetilde{\ell}_{n}(\Sigma)-U_{n}\le-b_{k}n+A_{1}D_{*}\sqrt{n\log\log n},
\]
where $D_{*}=\sup\{\|\Sigma-\Gamma\|_{F}:\Sigma,\Gamma\in\mathcal{K}\}<\infty$.
Because $\sqrt{n\log\log n}=o(n)$, the right-hand side is negative
for all sufficiently large $n$, almost surely. Taking suprema over
the near and far regions shows that 
\[
T_{k,n}\le U_{n}+\frac{A_{1}^{2}}{4c_{k}}\log\log n
\]
eventually almost surely. The lower bound $T_{k,n}\ge U_{n}$ is immediate
from $\mathcal{G}_{*}\subset\mathcal{F}_{k}$. 
\end{proof}

\end{document}
